\documentclass[10pt,a4paper]{amsart}
\usepackage[margin=19mm]{geometry}
\usepackage{amsmath,amssymb,mathtools}
\usepackage[scr=rsfs,cal=euler]{mathalfa}
\usepackage{xfrac}
\usepackage[unicode,hidelinks,bookmarks=false]{hyperref}
\newcommand{\ie}{i.e.\ }
\newcommand{\cf}{cf.\ }
\newcommand{\resp}{respectively\ }
\newcommand{\Subsection}{\subsection}
\providecommand{\nobreakdash}{\nobreak\hbox{-}}
\setlength{\emergencystretch}{2em}
\multlinegap0pt
\usepackage{amssymb}
\usepackage{mathtools}
\usepackage{tikz}
\usepackage[shortlabels]{enumitem}
\usepackage{xcolor}
\usepackage{dsfont}
\usepackage[normalem]{ulem}
\newtheorem{theorem}{Theorem}
\newtheorem{lemma}{Lemma}
\newtheorem{claim}{Claim}
\def\e{\epsilon}
\def\l{\ell}
\title{Recursive upper bounds for the vertex online Ramsey game with applications to hypergraph Ramsey numbers}
\author{D\'aniel Dob\'ak, Eion Mulrenin}
\date{Source: arXiv:2605.16607v1 (CC BY 4.0)}
\begin{document}
\maketitle

\noindent\emph{Source and licence.} Recursive upper bounds for the vertex online Ramsey game with applications to hypergraph Ramsey numbers. Version arXiv:2605.16607v1 (CC BY 4.0), distributed by arXiv under the Creative Commons Attribution 4.0 International licence, http://creativecommons.org/licenses/by/4.0/, as recorded in the arXiv licence metadata for this exact version and shown on the abstract page https://arxiv.org/abs/2605.16607v1. Full licence notice: \texttt{licenses/arxiv-2605.16607-CC-BY-4.0.txt}.\par\smallskip

\noindent\textit{Editorial scope.} Counted unit: the article's theorem bound (source0.tex lines 351--362). The article's complete original proof of it is delivered with it (source0.tex lines 364--507, one \texttt{proof} environment of 144 lines). No mathematics is written, repaired or re-derived: the statement and the proof are the article's own lines, in the article's own notation, and no retained line is substituted or re-flowed.  No further source-local context is needed: every cross-reference of every retained line resolves inside the two delivered spans.  Nothing was omitted from any retained span, so the package carries no omission record.  No retained line contains a citation, so the wrapper carries no bibliography and no bibliography entry.  No retained line contains a figure environment, an image-inclusion command or drawing code, so no image is delivered and the package declares no asset dependency.  Editorial formatting only, in addition: an AMS \texttt{amsart} preamble that restates the article's own declarations and macros used by the retained lines, namely the environments \texttt{theorem}, \texttt{lemma}, \texttt{claim} on separate counters, together with the standard packages those lines need.  The retained environments print in this document as Theorem 1, Lemma 1, Claim 1 in order of appearance, whereas the article prints the same items with the numbering of its own full document.  The complete native source of the article as distributed in the arXiv e-print for this exact version is bundled unchanged as \texttt{sources/200728/source0.tex}.
\begin{theorem}
    Let $2 \leq k < t_1, \dots, t_q$, and let $m = \tilde{r}_{k-1}(t_1-1, \dots, t_q-1)$.
    Then builder has a strategy in the $q$-color $k$-uniform online Ramsey game for cliques of size $t_1, \dots, t_q$ using at most $q^m+ k - 2$ vertices and at most $m \cdot q^m$ edges.
    In particular,
    \begin{equation} \label{eq: q-color ramsey bound}
        r_k(t_1, \dots, t_q) \leq q^m + k - 2,
    \end{equation}
    and
    \begin{equation} \label{eq: q-color online bound}
        \tilde{r}_k(t_1, \dots, t_q) \leq q^{m + \log_q m}.
    \end{equation}
\end{theorem}

\begin{proof}
Let $2 \leq k < t_1, \dots, t_q$ and $q \geq 2$ be given, and recall that we assume there is a builder strategy in the $(k-1)$-uniform $q$-color vertex online Ramsey game which forces a
monochromatic $K_{t_c-1}^{(k-1)}$ in some color $c \in [q]$ using
at most $m = \tilde{r}_{k-1}(t_1-1,\dots, t_q-1)$ edges against {\it any} painter strategy.
Our goal is to show that builder has a strategy in the $k$-uniform $q$-color vertex online Ramsey game which forces painter to draw a
monochromatic $K_{t_{c'}}^{(k)}$ in some color $c' \in [q]$
and which uses at most $q^m + k-2$ vertices and at most $m \cdot q^m$ edges.
Roughly speaking, the idea of the $k$-uniform builder strategy is that, once the $j$th vertex $v_j$ is revealed to begin the $j$th round, builder will run the $(k-1)$-uniform strategy as a ``sub-routine" on a certain set $R(v_j)$ of previously revealed vertices.
On this set, builder will run this optimal $(k-1)$-uniform strategy against an auxiliary $(k-1)$-uniform online coloring $\Delta_j$ which is determined by painter's $k$-uniform coloring $\Delta$.
By design, once there is a monochromatic $K_{t_c-1}^{(k-1)}$ under one of these auxiliary colorings $\Delta_j$, there will also be a monochromatic $K_{t_c}^{(k)}$ under $\Delta$.

More formally, we define builder's $k$-uniform strategy by recursively constructing a map $\Phi$ from the ordered vertex set $V = \{v_1, v_2, \dots\}$ into 
$[q]^*$, the set of all finite strings over the alphabet $[q]$.
For each $j \geq 1$, the $j$th round of the game begins with the vertex $v_j$ being revealed, and after playing this round, builder will define $\Phi(v_j) \in [q]^*$ and an associated set of previously revealed vertices $R(v_j) \subseteq \{v_1, \dots, v_{j-1}\}$.
$\Phi$ will be an injection on $V \setminus \{v_1, \dots, v_{k-2}\}$, and $R(v_j)$ will be the vertex set on which the $(k-1)$-uniform coloring $\Delta_j$ is defined.

To initialize, let $v_1, ..., v_{k-1}$ be the first $k-1$ vertices revealed in the game, and set
$\Phi(v_j) := \e$, the empty string, for each $j \in [k-1]$.
Now let $j \geq k$ and let $v_j$ be the vertex revealed at the beginning of the $j$th round, and recursively assume that builder has already defined $\Phi(v_i)$ and $R(v_i)$ for all $i < j$.
We define $\Phi(v_j)$ and $R(v_j)$ by the following recursive sub-procedure:
\begin{itemize}
    \item {\bf Initial step.} Set $\Phi_0(v_j) := \e$ and $R_0(v_j) := \emptyset$.
    \item {\bf Recursion step.} Suppose for $f \geq 0$ that builder has already defined $\Phi_f(v_j) \in [q]^*$ and a set of previously revealed vertices $R_f(v_j) \subseteq \{v_1, \dots, v_{j-1}\}$ of size $f$.
    Furthermore, suppose that builder has run $f$ rounds of the $(k-1)$-uniform strategy on this set $R_f(v_j)$ against an auxiliary $(k-1)$-uniform coloring $\Delta_j$ that she has defined, in the order that vertices were added to $R_f(v_j)$.
    
    If there exists a vertex $v_i \notin R_f(v_j)$ with $i < j$ such that $\Phi_f(v_j) = \Phi(v_i)$,
    %and $R_f(v_i) = R(v_j)$, 
    builder picks $i$ which is minimal with this property\footnote{Strictly speaking, $\Phi$ will be an injection on $V \setminus \{v_1, \dots, v_{k-2}\}$, so this is only necessary for $f \in [k-2]$.}
    and initiates round $f+1$ of the $(k-1)$-uniform vertex online Ramsey game on $R_f(v_j) \cup \{v_i\}$ with $v_i$ as the revealed vertex.
    At each step of the $(f+1)$st round, by the rules of the game, the $(k-1)$-uniform strategy will build an edge $e \cup \{v_i\}$, where $e \subseteq R_f(v_j)^{(k-2)}$, and ask for it to be colored immediately by the auxiliary coloring $\Delta_j$.
    Builder decides the color $\Delta_j(e \cup \{v_i\})$ by first, in the $k$-uniform game, revealing the edge $e \cup \{v_i, v_j\}$ (recall that we are in the $j$th round of the $k$-uniform game) and then having painter assign it a color $\Delta(e \cup \{v_i,v_j\}) \in [q]$ immediately.
    Once this has occurred, builder sets
    \begin{equation*}
        \Delta_j(e \cup \{v_i\}) := \Delta(e \cup \{v_i, v_j\}).
    \end{equation*}
    Now that the edge $e \cup \{v_i\}$ has received a color under $\Delta_j$ in the $(k-1)$-uniform game, we proceed to the next step of round $f+1$ of this game.

    At the conclusion of the $(f+1)$st round, builder sets:
    \begin{itemize}
        \item $R_{f+1}(v_j) := R_f(v_j) \cup \{v_i\}$; and
        \item $\Phi_{f+1}(v_j)$ to be the string in $[q]^*$ obtained by appending, in order, the colors $\Delta_j(e \cup \{v_i\})$ used at each step in the $(f+1)$st round of the $(k-1)$-uniform game to the word $\Phi_f(v_j)$.
    \end{itemize}

    Eventually, for some $f$, there will be no $i < j$ with $\Phi(v_i) = \Phi_f(v_j)$.
    When this happens, builder sets $\Phi(v_j) := \Phi_f(v_j)$ and $R(v_j) := R_f(v_j)$, and ends round $j$ of the $k$-uniform game.
\end{itemize}

\begin{lemma}
\label{lem: main}
    Let $k \leq i < j$, and suppose $v_i \in R(v_j)$.
    Then the following claims hold:
    \begin{itemize}
        \item[(i)] $R(v_i) = \{v_\l \in R(v_j): \l < i\}$; and
        \item[(ii)] If $e \subseteq R(v_i) \cap R(v_j) \overset{(i)}{=} R(v_i)$ has size $k-1$, then builder drew the edge $e \cup \{v_i\}$ if and only if she drew the edge $e \cup \{v_j\}$, and in this event,
        \begin{equation*}
            \Delta(e \cup \{v_i\}) = \Delta(e \cup \{v_j\}). 
        \end{equation*}
    \end{itemize}
\end{lemma}

\begin{proof}
    {\it (i)}
    By construction, $v_i \in R(v_j)$ if and only if there exists $f \geq 0$ with $\Phi_f(v_j) = \Phi(v_i)$, so it's enough to prove that for this same choice of $f$, $R_f(v_j) = R(v_i)$.
    Write $\Phi_f(v_j) = \Phi(v_i)$ as both $w_1 \cdots w_f$ and $w_1' \cdots w_g'$, where $w_h$ and $w_h'$ are, respectively, the substrings appended after round $h$ of the $(k-1)$-uniform game when builder was defining $\Phi_f(v_j)$ and $\Phi(v_i)$.
    Note that since $\Phi$ is an injection on $V \setminus \{v_1, \dots, v_{k-2}\}$, when $k \leq i < j$ we can write
    \begin{equation}
    \label{eq: Rs}
        R_f(v_j) = \Phi^{-1}(w_1) \cup \dots \cup \Phi^{-1}(w_1 \cdots w_f)
        \hspace{0.25cm}
        \text{and}
        \hspace{0.25cm}
        R(v_i) = \Phi^{-1}(w_1') \cup \dots \cup \Phi^{-1}(w_1' \cdots w_g').
    \end{equation}
    We now prove by induction that $f = g$ and $w_h = w_h'$ for all $h \in [f]$, which suffices to prove $R_f(v_j) = R(v_i)$ by~\eqref{eq: Rs}.
    %To see that this suffices to prove the claim, observe that since builder's $(k-1)$-uniform strategy was deterministic, and $\Phi$ was specifically constructed as an injection after the first $k-2$ vertices, the same vertex $u_g$ was necessarily added to both $R_{g-1}(v_j) \subseteq R_f(v_j)$ and $R_{g-1}(v_i) \subseteq R(v_i)$ once the word $w_g$ was appended to $\Phi_{g-1}(v_j)$ and $\Phi_{g-1}(v_i)$.

    For the base case, note that $w_h = w_h' = \e$ for all $h \in [k-2]$, as builder's $(k-1)$-uniform strategy cannot build any edges until there are at least $k-1$ vertices revealed.
    Now suppose that $w_1 = w_1'$, \dots, $w_h = w_h'$ for some $h$ with $k-2 \leq h < \min\{f,\l\}$, and that moreover, in the first $h$ rounds, builder's $(k-1)$-uniform strategy picked the same edges in the same order against both $\Delta_j$ and $\Delta_i$, and that each such edge was assigned the same color by $\Delta_j$ and $\Delta_i$.
    At this stage, there must have been some vertex $v_\l$ with $k-1 \leq \l < i, j$ such that $\Phi(v_\l) = w_1 \cdots w_h = w_1' \cdots w_h'$, and this $v_\l$ is unique, as again, $\Phi$ is an injection on vertices past the first $k-2$.
    In both $(k-1)$-uniform games, this means $v_\l$ was the vertex revealed at the beginning of the $(h+1)$st round, and $w_{h+1}$ and $w_{h+1}'$ are the sequences of colors that were used by $\Delta_j$ and $\Delta_i$, respectively, on edges of the form $e \cup \{v_\l\}$.

    Since the first $h$ rounds of the $(k-1)$-uniform game against $\Delta_j$ and $\Delta_i$ were identical, when $v_\l$ is revealed to begin the $(h+1)$st round, the deterministic $(k-1)$-uniform strategy will build the same edges, and since $\Phi_f(v_j) = \Phi(v_i)$, it must be that each time a new edge is built, it receives the same color from $\Delta_j$ and $\Delta_i$.
    And again, by the fact that the $(k-1)$-uniform strategy is deterministic, the $(h+1)$st round will conclude in both games at exactly the same step, whence $w_{h+1}=w_{h+1}'$ and the $(h+1)$st round is identical in both games.

    {\it (ii)} now follows immediately from the fact that the first $f$ rounds of both $(k-1)$-uniform games are identical.
    %The first point follows from the proof above: since $\Phi_f(v_j)$ and $\Phi(v_i)$ are the same, builder has necessarily picked the same sequence of edges prior to $e$, and they were colored identically by $\Delta_i$ and $\Delta_j$.
    %Thus, once builder comes to $e$, their $(k-1)$-uniform strategy takes $e$ against $\Delta_i$ if and only if it takes $e$ against $\Delta_j$, and then again using $\Phi_f(v_j) = \Phi(v_i)$, it must be that $\Delta_i(e) = \Delta_j(e)$, i.e., 
    %\begin{equation*}
    %    \Delta(e \cup \{v_i\}) = \Delta(e \cup \{v_j\}). 
    %\end{equation*}
\end{proof}

\begin{claim}
\label{clm: endgame}
    If there is a vertex $v_n$ with $\Phi(v_n)$ of length $m$, then painter has drawn a monochromatic $K_{t_c}^{(k)}$ in some color $c \in [q]$ under $\Delta$, and thus the game ends.
\end{claim}

\begin{proof}
    First, note that a label of length $m$ implies that builder's $(k-1)$-uniform strategy built $m$ $(k-1)$-uniform edges in $R(v_n)$, which were colored under $\Delta_n$, and by assumption, this strategy guaranteed a monochromatic $K_{t_c-1}^{(k-1)}$ in some color $c \in [q]$ once $m$ edges were drawn.
    Let $t=t_c$ and let $U = \{u_1, \dots, u_{t-1}\} \subseteq R(v_n)$ be the vertices of this $K_{t-1}^{(k-1)}$ which is monochromatic under $\Delta_n$.
    Our claim is that the set $U \cup \{v_n\}$ induces a monochromatic $K_t^{(k)}$ in color $c$ under $\Delta$.
    First, note that 
    \begin{equation*}
        \Delta(e \cup \{v_n\}) = c
    \end{equation*}
    for all $e \in U^{(k-1)}$;
    indeed, this follows since each such edge $e$ was colored with $c$ by $\Delta_n$.

    Now let $i_1 < \dots < i_k \in [t-1]$.
    Since $u_{i_k} \in R(v_n)$, Lemma~\ref{lem: main}(i) implies that
    $R(u_{i_k}) \supseteq \{u_{i_1}, \dots, u_{i_{k-1}}\}$, and so by Lemma~\ref{lem: main}(ii), since builder drew the edge $\{u_{i_1}, \dots, u_{i_{k-1}}, v_n\}$, she also drew the edge $\{u_{i_1}, \dots, u_{i_k}\}$, and we have
    \begin{equation*}
        \Delta(u_{i_1}, \dots, u_{i_{k-1}}, u_{i_k}) = \Delta(u_{i_1}, \dots, u_{i_{k-1}}, v_n) = c,
    \end{equation*}
    as promised.
\end{proof}

\begin{claim}
\label{clm: injection}
    Let $V = \{v_1, \dots, v_n\}$ be the set of all vertices revealed in the game.
    Then $\Phi$ is an injection from $V \setminus \{v_1, \dots, v_{k-2}, v_n\}$ into $[q]^{\leq m-1}$.
\end{claim}

\begin{proof}
    That $\Phi$ is an injection on this set of vertices is immediate from its construction.
    The claim that the codomain is $[q]^{\leq m-1}$ follows from Claim~\ref{clm: endgame}, since the game ends as soon as vertex $v_n$ receives a label $\Phi(v_n)$ with at least $m$ letters.
\end{proof}

To finish, we just have to upper-bound the total number of vertices and edges we have used.
By Claim~\ref{clm: injection}, we immediately have
\begin{equation*}
    |V| \leq k-1+\sum_{i=0}^{m-1} q^i \leq q^m + k - 2.
\end{equation*}

For the edge count, observe that edges with largest vertex $v_j$ are in one-to-one correspondence with letters in the string $\Phi(v_j)$.
Indeed, if an edge $e = \{v_{i_1}, \dots, v_{i_k}\}$ was built by builder with $i_1 < \dots < i_k$ and painter assigned $\Delta(e) = c$, then the letter $c$ appears in the word $\Phi(v_{i_k})$ at the point when builder's $(k-1)$-uniform strategy picked the edge $e \setminus \{v_{i_k}\}$ during round $i_k$ of the $k$-uniform game.
Thus, the number of edges built by builder in the $k$-uniform game is exactly
\begin{eqnarray*}
    \sum_{j=1}^n |\Phi(v_j)|
    &\leq& m + \sum_{i=1}^{m-1} i \cdot q^i\\
    &\leq& m \cdot q^m,
\end{eqnarray*}
as required.
\end{proof}

\end{document}
